\section{Verified Multi-Word Compare-and-Set Algorithm}
\label{sec:verification}

The correctness of \Kcas transactions relies on that of the underlying MCAS algorithm (see \cref{sec:implementation_mcas}), whose implementation (\cref{fig:mcas_code_1,fig:mcas_code_2,fig:mcas_code_3}) is short but extremely subtle.
We verified the algorithm~\refLibTheoriesFile{zoo_mcas/mcas} using the \Zoo framework~\citep{DBLP:journals/pacmpl/AllainS26}, which provides (1) a translator from \OCaml to \Rocq and (2) a precise semantics for concurrent primitives (such as \ocamlinline{Atomic.compare_and_set}) and physical equality.
All our proofs are mechanized in the \Rocq proof assistant.

Note that \Zoo assumes a \emph{sequentially consistent} memory model whereas \OCamlFive has a \emph{relaxed} memory model~\citep{DBLP:conf/pldi/DolanSM18}.
We inherit the limitation and conjecture that this assumption does not invalidate our concurrent invariant, the core of the proof, and that it can be adapted to relaxed memory following the ``recipe'' of \citet{DBLP:journals/pacmpl/MevelJ21}.

\subsection{Specification}

\begin{figure}
\begin{mathpar}
    \iSpec[
      lab=make-spec
    ]{
      \iTrue
    }{
      \c{make}\ \zooVal
    }[
      \zooLoc
    ]{
      \locInv\ \zooLoc\ \iIname \iSep \\
      \zooLoc \iPointstoTwo \zooVal
    }
  \and
    \iAspec[
      lab=get-spec,
      mask=\iIname
    ]{
      \locInv\ \zooLoc\ \iIname
    }[
      \zooVal
    ]{
      \zooLoc \iPointstoTwo \zooVal
    }{
      \c{get}\ \zooLoc
    }{
      \zooLoc \iPointstoTwo \zooVal
    }[
      \zooValTwo
    ]{
      \zooVal \similar \zooValTwo
    }
  \and
    \iAspec[
      lab=cas-spec,
      mask=\iIname
    ]{
      \iBigSep_{\zooLoc \in \zooLoc{}s}
        \locInv\ \zooLoc\ \iIname
    }[
      \zooVals
    ]{
      \iBigSep_{\zooLoc, \zooVal \in \zooLoc{}s, \zooVals}
        \zooLoc \iPointstoTwo \zooVal
    }{
      \c{cas}\ \zooLoc{}s\ \v{befores}\ \v{afters}
    }[
      \zooBool
    ]{
      \myif{
        \zooBool
      }{
        \begin{inbox}
          \zooVals \similar \v{befores} \iSep \\
          \iBigSep_{\zooLoc, \zooVal \in \zooLoc{}s, \v{afters}}
            \zooLoc \iPointstoTwo \zooVal
        \end{inbox}
      }{
        \begin{inbox}
          \Exists i\ \zooVal\ \v{before}. \\
          \mapLookup{\zooVals}{i} = \Some{\zooVal} \iSep \\
          \mapLookup{\v{befores}}{i} = \Some{\v{before}} \iSep \\
          \zooVal \nonsimilar \v{before} \iSep \\
          \iBigSep_{\zooLoc, \zooVal \in \zooLoc{}s, \zooVals}
            \zooLoc \iPointstoTwo \zooVal
        \end{inbox}
      }
    }[
      \v{res}
    ]{
      \v{res} = \zooBool
    }
\end{mathpar}
\caption{MCAS: Specification}
\label{fig:mcas_spec}
\end{figure}


The specification of MCAS is given in \cref{fig:mcas_spec}.
The persistent assertion $\locInv\ \zooLoc\ \iIname$ represents the knowledge that \zooLoc is a valid MCAS location.
The exclusive assertion $\zooLoc \iPointstoTwo \zooVal$ represents the ownership of location \zooLoc and the knowledge that it contains value \zooVal.
$\c{make}\ \zooVal$ creates a new location initialized with \zooVal.
$\c{get}\ \zooLoc$ atomically reads the content of \zooLoc.
$\c{cas}\ \zooLoc{}s\ \v{befores}\ \v{afters}$ performs an MCAS operation on locations $\zooLoc{}s$ with expected values \v{befores} and desired values \v{afters}; on success, it atomically updates $\zooLoc{}s$ from \v{befores} to \v{afters}; on failure, it must have observed a location whose value was different than expected.

\subsection{Linearization}

\paragraph{External linearization.}

As explained in \cref{sec:implementation_mcas}, an MCAS operation may be helped and in fact linearized by another MCAS.
Consequently, the linearization point of both \c{get} and \c{cas} may be external.
In fact, a helper may linearize many operations at once, including the original MCAS and other helpers.
Also, a helper may linearize a helped MCAS while helping yet another MCAS.

\paragraph{Future-dependent linearization.}

When two operations helping an MCAS detect an unexpected value, they both attempt to atomically update its status to B.
The question is: which of them linearizes the MCAS operation (and other helpers) at the point when it observes an unexpected value?
The answer is: the one which ``wins'' the update.
Consequently, the linearization point of \c{cas} may also be future-dependent.

\subsection{Challenges}

% + prophecy variables (Zoo.proph, Zoo.resolve, Zoo.resolve_with)
% + ghost identifiers (Zoo.id)

The implementation is fairly short (one hundred lines of code) but extremely subtle.
We outline the main challenges.

\paragraph{Mutually recursive invariants.}

Due to the cyclic structure of ongoing MCASes (see \cref{fig:loc_representation_3}), the definitions of the location invariant (\locInv) and the MCAS invariant are mutually recursive.
Thankfully, \Iris provides a way to define such fixpoints.

Another difficulty comes from the fact that the creation of a location involves a dummy MCAS operation whose invariant has to be initialized at the same time as the location's.
We had to introduce a way to break the cycle.

\paragraph{Logical state.}

The verification involves a monotonic logical status that is not determined by the physical status: the physical U status is divided into indexed logical U status indicating the progress made by the MCAS operation.

\paragraph{External linearization.}

To handle external linearization of an MCAS operation and its helpers by a helper, their atomic updates%
\footnote{
In \Iris, an \emph{atomic update} is a resource that materializes a linearization point.
\Citet{DBLP:journals/pacmpl/MulderK23} provide a good description.
}
are stored into the MCAS invariant.
At the linearization point, the winning helper triggers all the atomic updates and puts them back into the invariant to be retrieved later by their respective owner.

\paragraph{Global prophecy variable.}

To handle future-dependent linearization, we predict both the winner and the outcome of an MCAS operation through a shared prophecy variable.
To distinguish the winner, we use the same trick as \citet{DBLP:journals/pacmpl/JungLPRTDJ20} in their proof of RCDSS: each MCAS participant is assigned a unique physical identifier (implemented using a prophecy variable) that is part of the prediction.

\paragraph{Local prophecy variable.}

For subtle reasons related to the semantics of physical equality~\citep{DBLP:journals/pacmpl/AllainS26}, we also need to introduce a local prophecy variable in \c{cas}.
Essentially, the problem is that we cannot predict the outcome of physical equality in advance just by looking at the values; in other words, equality in \Rocq does not imply physical equality in \OCaml.

% generative constructors

\paragraph{MCAS history.}

Another subtle point in the algorithm is the fact that a location cannot be locked more than once by a single MCAS operation.
While this is obvious when reading the code, it is not in the proof.
To enforce it logically, each location maintains a ghost history of distinct MCASes; all MCAS in this history are finished except the more recent one.
